\section{题07：核心任务调度（P16789）}

\subsection{问题重述}

有 $N$ 个任务，每个任务恰好需要一天完成，小蓝每天至多完成一个任务，也可以在某天不安排任何任务。天数从第 $1$ 天开始按正整数编号。每个任务 $i$ 有一个截止时间 $t_i$ 和一个价值 $w_i$：如果在第 $t_i$ 天或之前完成该任务，则获得价值 $w_i$；若超过截止时间，任务失效且不可再完成。此外，部分任务被标记为核心任务（$\text{is\_key}_i = 1$），其余为普通任务（$\text{is\_key}_i = 0$）。公司要求小蓝制定调度方案时，必须先保证完成的核心任务数量尽可能多，在所有满足这一要求的方案中再使总价值尽可能大。求最多能完成的核心任务数 $K$ 和该前提下的最大总价值 $V$。约束：$N \le 2 \times 10^5$，$t_i \le 10^9$，$w_i \le 10^9$。

样例中共 $4$ 个任务：任务 $1$ 的 $t=3, w=10$ 且为核心任务，任务 $2$ 的 $t=1, w=5$ 且为核心任务，任务 $3$ 的 $t=2, w=100$ 且为普通任务，任务 $4$ 的 $t=1, w=100$ 且为普通任务。最优方案为第 $1$ 天完成任务 $2$、第 $2$ 天完成任务 $3$、第 $3$ 天完成任务 $1$，完成 $2$ 个核心任务，总价值 $5+100+10=115$。

\subsection{从蛮力枚举到带权调度贪心}

先考虑最朴素的思路：枚举所有任务的子集，检查该子集是否存在一个可行调度（即能否将每个选中任务分配到一个不超过其截止时间且互不冲突的日期），若可行则记录其核心任务数和总价值，最终在所有可行子集中取字典序最优者。可行性检查本身就是一个经典问题：按截止时间排序后贪心分配最早可用日期，若所有选中任务都能分配到日期则可行。枚举 $2^N$ 个子集、每个子集花费 $\mathrm{O}(N \log N)$ 检查，总复杂度 $\mathrm{O}(2^N N \log N)$，在 $N = 2 \times 10^5$ 面前毫无希望。

但这道题的底层结构正是经典的带截止时间单位任务调度问题：$N$ 个任务各耗时一单位，各有截止时间和价值，要求选出一个可调度子集使总价值最大。该问题有一个优雅的拟阵贪心算法。将所有任务按截止时间从小到大排序，依次处理每个任务：将其价值加入一个小根堆，若此时堆的大小（即已选任务数）超过了当前任务的截止时间，说明时间槽不够用了，则弹出堆中价值最小的任务。处理完所有任务后，堆中剩余的即为所选任务。

以题目样例演示该算法（暂时忽略核心/普通的区分，仅按原始价值 $w_i$ 运行）。四个任务按截止时间排序后为：$(t=1,w=5)$、$(t=1,w=100)$、$(t=2,w=100)$、$(t=3,w=10)$。依次处理：第一个任务入堆，堆大小 $1 \le t=1$，保留；第二个任务入堆，堆大小 $2 > t=1$，弹出最小值 $5$，堆中只剩价值 $100$ 的任务；第三个任务入堆，堆大小 $2 \le t=2$，保留；第四个任务入堆，堆大小 $3 \le t=3$，保留。最终堆中有三个任务（价值 $100,100,10$），总价值 $210$。然而正确的最优答案是完成两个核心任务（价值 $5$ 和 $10$）加上普通任务 $3$（价值 $100$），总价值 $115$。可见直接最大化总价值的贪心无法保证核心任务数量优先。

\subsection{字典序双目标的数值编码}

本题的真正目标是一个字典序优化：第一关键字是核心任务数量 $K$（最大化），第二关键字是总价值 $V$（最大化）。等价地，我们希望找到一个方案使得二元组 $(K, V)$ 在字典序下最大。这提示我们可以将双目标编码为单一数值，从而复用经典的单目标贪心框架。

关键观察：设所有任务的原始价值之和为 $S = \sum_{i=1}^{N} w_i$，令 $C = S + 1$。定义每个任务的调整价值为
\[
\tilde{w}_i = w_i + C \cdot \text{is\_key}_i
\]
那么，对任意两个方案 $A$ 和 $B$，$\sum_{i \in A} \tilde{w}_i > \sum_{i \in B} \tilde{w}_i$ 当且仅当 $A$ 的字典序 $(K_A, V_A)$ 严格优于 $B$ 的字典序 $(K_B, V_B)$。原因如下：任何多完成一个核心任务的方案，其调整价值之和至少多出 $C$（因为新完成的核心任务贡献了 $C + w_i \ge C$）；而任何不改变核心任务数量的方案，其调整价值之和的差异至多为 $\sum w_i = S < C$。因此调整价值之和的比较天然地实现了先比较 $K$、再比较 $V$ 的字典序语义。

在样例中，$S = 10 + 5 + 100 + 100 = 215$，$C = 216$。四个任务的调整价值分别为 $\tilde{w}_1 = 10 + 216 = 226$，$\tilde{w}_2 = 5 + 216 = 221$，$\tilde{w}_3 = 100 + 0 = 100$，$\tilde{w}_4 = 100 + 0 = 100$。现在以调整价值 $\tilde{w}_i$ 替代原始价值 $w_i$ 运行前述带权调度贪心算法。

\subsection{贪心算法的逐步推演}

将四个任务按截止时间排序（截止时间相同时任意顺序均可，因为堆会在超出截止时间时自动淘汰价值最低者）：依次为 $(t=1, \tilde{w}=221)$、$(t=1, \tilde{w}=100)$、$(t=2, \tilde{w}=100)$、$(t=3, \tilde{w}=226)$。

第一步：处理截止时间 $1$、调整价值 $221$ 的任务 $2$。将其入堆，堆大小 $1 \le 1$，保留。此时堆中为 $\{221\}$。

第二步：处理截止时间 $1$、调整价值 $100$ 的任务 $4$。将其入堆，堆大小变为 $2$，超过了当前截止时间 $1$，触发淘汰。弹出堆中最小调整价值 $100$，堆中恢复为 $\{221\}$。

第三步：处理截止时间 $2$、调整价值 $100$ 的任务 $3$。将其入堆，堆大小 $2 \le 2$，保留。堆中为 $\{221, 100\}$。

第四步：处理截止时间 $3$、调整价值 $226$ 的任务 $1$。将其入堆，堆大小 $3 \le 3$，保留。堆中最终为 $\{221, 100, 226\}$，对应任务 $2$、$3$、$1$。

现在从最终堆中还原答案。核心任务数 $K$ 等于最终堆中核心任务的出现次数，在此为 $2$（任务 $2$ 和任务 $1$）。总价值 $V$ 等于最终堆中所有任务的原始价值之和，在此为 $5 + 100 + 10 = 115$。这与题目给出的最优解完全一致。

\subsection{为什么编码技巧是正确的}

上述编码技巧的正确性依赖于 $C > \sum w_i$ 这一不等式。直观上，$C$ 充当了“核心任务”这一维度的权重，使得任何核心任务数量的差异都能压倒总价值维度的任何差异。形式化地，考虑任意两个可行方案 $A$ 和 $B$，它们的调整价值总和之差为
\[
\Delta = \sum_{i \in A} \tilde{w}_i - \sum_{i \in B} \tilde{w}_i = C \cdot (K_A - K_B) + (V_A - V_B).
\]
若 $K_A > K_B$，则 $K_A - K_B \ge 1$，故 $\Delta \ge C - (V_B - V_A) \ge C - S = 1 > 0$，方案 $A$ 的调整价值更大。若 $K_A = K_B$，则 $\Delta = V_A - V_B$，调整价值的比较退化为纯价值比较。因此调整价值之和的最大化恰好等价于字典序 $(K, V)$ 的最大化。

值得注意的是，这一步编码并不改变贪心算法本身的正确性——带权调度贪心的正确性根植于该问题所对应的拟阵结构（独立集为所有存在可行调度的任务子集），调整价值仅是改变了每个元素的权重，拟阵的贪心算法对任意非负权重均成立。

\subsection{复杂度与实现细节}

整个算法分为三个阶段：计算权重和并构造调整价值，按截止时间排序，以及单次扫描加堆操作。第一阶段遍历所有任务一次，$\mathrm{O}(N)$；排序阶段 $\mathrm{O}(N \log N)$；扫描阶段每个任务至多入堆一次、出堆一次，堆操作 $\mathrm{O}(\log N)$，故 $\mathrm{O}(N \log N)$。总时间复杂度 $\mathrm{O}(N \log N)$，空间复杂度 $\mathrm{O}(N)$。在 $N = 2 \times 10^5$ 的约束下，$\log_2 N \approx 18$，总操作量约 $3.6 \times 10^6$ 级别，远在时限之内。

实现时有几个值得注意的细节。由于 $w_i$ 可达 $10^9$ 且 $N$ 可达 $2 \times 10^5$，$S = \sum w_i$ 可达 $2 \times 10^{14}$，需使用 \mil{int64_t}（即 \mil{long long}）存储 $C$ 和调整价值，避免溢出。排序时直接对截止时间排序即可，因为贪心正确性不依赖同截止时间内任务的处理顺序——堆会在总任务数超出截止时间时自动淘汰最劣者，无论它们以何种顺序入堆。另外，在扫描过程中，与其从最终的调整价值反推核心任务数和总价值，不如在入堆和出堆时直接维护 \mil{total_val} 和 \mil{core_cnt} 两个变量——入堆时加上对应任务的 $w_i$ 和 $\text{is\_key}_i$，出堆时减去——这样最终答案可以直接输出，无需额外解码步骤。堆中存储的是（调整价值，原始下标）的 pair，因为出堆时需要知道被淘汰的是哪个任务，以便正确更新 \mil{total_val} 和 \mil{core_cnt}。

\subsection{代码实现}
\inputminted{c++}{../src/p07_task_scheduler.cpp}
